\section{Discussion}
\label{sec:discussion}

The proposed framework establishes a clear separation of concerns among its core components: local information weights quantify supervised distributional associations, granular dependency evaluates training representation coverage, and the inner-fold loss guard delineates a predictive feasibility boundary. A primary theoretical insight of this decoupling is that maximizing training-internal representation quality does not universally guarantee out-of-sample prediction preservation. By acknowledging non-monotone granular gains, the search mechanism naturally accommodates structural reformations in granular space without compromising exact cost-based branch pruning.

Our empirical findings provide critical insights into the operational trade-offs of cost-sensitive LDL reduction. Across twenty-five native benchmarks, the framework identifies highly informative attribute subsets at ultra-low acquisition costs compared to matched label-free baselines. Crucially, the multi-domain expansion reveals that the near-degenerate regime first observed on \emph{Yeast-cold} and \emph{Human Gene} is the dominant case rather than an exception: on twenty of the twenty-five benchmarks, the mean-distribution predictor itself achieves optimal audit performance under the Euclidean $k$-NN protocol. In such cases, returning an empty set at zero cost is not a predictive failure, but rather the theoretically correct and most cost-effective response. Furthermore, scenarios where nonempty subsets underperform the empty baseline (e.g., boundary cases under uninformative predictors) highlight the essential role of internal feasibility checks in preventing unnecessary feature acquisition. While granular guidance effectively navigates complex search spaces to discover non-dominated low-cost candidates, its performance gains interact with dataset geometry and cost distribution topologies.

To maintain methodological rigor, future confirmatory evaluations should preserve the frozen algorithmic structure while expanding the experimental scope across diverse multi-source cost scenarios and heterogeneous LDL predictors. In addition, assessing exact search methods on compact real-world instances will further illuminate the trade-offs between heuristic optimality gaps and computational overhead.

\section{Conclusion}
\label{sec:conclusion}

In this paper, we presented a novel cost-constrained attribute reduction framework tailored for Label Distribution Learning (LDL). By combining local distribution-information weighting with a source-dependent directed dissimilarity, our approach captures both supervised label associations and continuous geometric structures without forcing rigid metric symmetry. We established that distance monotonicity under fixed statistics does not automatically guarantee monotone granular dependency—a theoretical insight that motivates our non-monotone granular-ball representation. To address this, we cast attribute reduction as a minimum-acquisition-cost optimization task subject to an internal prediction-loss constraint, decoupling dynamic granular dependency guidance from predictive feasibility guarantees. 

Extensive experiments across twenty-five native LDL benchmarks demonstrate that the proposed framework cuts feature acquisition costs by over two orders of magnitude while strictly maintaining predictive feasibility. Benchmark comparisons against five published scoring cores spanning 2022--2025 confirm that our method avoids exorbitant acquisition expenses for equivalent or superior prediction loss. Crucially, our empirical analysis uncovers near-degenerate regimes where the zero-cost baseline is optimal, underscoring the necessity of evaluating internal feasibility rather than enforcing artificial "lossless" reduction; the guard ablation shows that guard-free selection actively wastes budget on these domains, and the predictor-sensitivity study shows the same guard certifying low-cost dominant subsets on revived domains once a suitable predictor is deployed, on six of nine such datasets outperforming both the zero-cost and full-cost references at $1$--$6\%$ of full-set cost.

Future research will extend this framework to online and streaming LDL scenarios, investigate adaptive prediction-loss bounds under multi-source cost heterogeneity, and explore exact constraint-solving techniques for large-scale real-world deployments.
